\documentclass{article}
\usepackage{graphicx, hyperref, booktabs, listings, amsmath}
\title{Vyasa: Structured Architectural Graphs with\\Deterministic Witness for Prompt-to-Code Generation}
\author{Prashant Pandey}
\date{July 2026}

\begin{document}
\maketitle

\begin{abstract}
Current AI code generators produce applications token-by-token from natural language prompts, resulting in cross-file inconsistencies, hallucinated imports, and missing foreign keys. Production survival rates for prompt-to-code tools range from 38\% (Lovable) to 71\% (v0), with median refactor costs of 2.4--4.8$\times$ the original generation time. We present \textbf{Vyasa}, a three-stage compiler architecture where the LLM produces a typed \textbf{Architectural Graph} (structured intermediate representation), a deterministic \textbf{Witness} validates it against 12 structural rules, and a \textbf{Manifest Engine} compiles validated graphs into working code with zero hallucination. The Architectural Graph compresses application specification approximately 13:1 versus direct token generation. We train a 3B-parameter adapter (Qwen-2.5-3B-Instruct + QLoRA) on teacher-generated graphs, achieving 66.7\% valid JSON with 25.0\% Witness-clean rate. Deterministic stages execute in under 2 milliseconds total. We introduce \textbf{GRPO+Witness}, a training flywheel where the Witness itself provides the reward signal---no human evaluation is needed. The deterministic compilation guarantees that every import resolves, every type is consistent, and every foreign key exists, regardless of the LLM's underlying error rate.
\end{abstract}

\section{Introduction}
\subsection{The Token-by-Token Problem}
Current LLM-based code generation tools---Bolt.new, Lovable, v0, Replit Agent---generate code autoregressively, one token at a time. The consequences: cross-file type inconsistency, hallucinated imports, missing foreign keys, and compounding errors.

A 2026 audit of 31 AI-generated codebases found that prompt-to-code tools produce applications with 38--50\% production survival rates after 60 days~\cite{appycodes}. The median refactor cost multiplier---hours of cleanup per hour of prompting---ranges from 2.4$\times$ to 4.8$\times$.

\subsection{The Structured IR Alternative}
We propose: \textbf{the LLM should not write code. It should write architecture. A deterministic compiler should write code.}

Two concurrent works arrived at similar architectures. \textbf{PlanCompiler}~\cite{plancompiler} achieves 96\% first-pass success using GPT-4 with typed node registry and static graph validation. \textbf{Compiled AI}~\cite{compiledai} achieves 57$\times$ token reduction via compile-time generation and zero-token execution.

\textbf{Vyasa differs:} (1) we train a local 3B model rather than relying on cloud APIs, (2) we introduce the Witness as a training reward signal (GRPO+Witness), and (3) we apply the structured IR approach to both code generation AND conversation.

\section{Architecture}
Vyasa is a three-stage compiler: Conception (LLM produces Architectural Graph) $\rightarrow$ Witness (12 deterministic structural checks) $\rightarrow$ Manifest (compiler produces code files).

The Architectural Graph is a typed JSON IR describing entities, relations, routes, pages, components, and auth. It compresses application specification $\sim$13:1 versus direct token generation.

When the Witness finds issues, the system enters a Prana-Apana refinement loop (up to 8 iterations) where the LLM receives issue feedback and regenerates.

\section{Training}
We use a teacher-student pipeline: DeepSeek-Chat $\rightarrow$ Witness filter $\rightarrow$ clean pairs only. The adapter is Qwen-2.5-3B-Instruct-4bit + QLoRA (rank 8, 8 layers, 300 iterations).

\textbf{GRPO+Witness:} The Witness IS the reward function. No human evaluation. For G=4 rollouts per prompt, $A_i = (R_i - \text{median}(R)) / \text{std}(R)$. Loss = $-\text{mean}(A_i \cdot \log \pi(a_i|s))$. The Witness never changes; the model converges toward it.

\section{Evaluation}
We evaluate on a 50-description test suite spanning 8 domains.

\begin{table}[h]
\centering
\begin{tabular}{lcccc}
\toprule
Model & Valid JSON & Witness Clean & Latency & Graph Size \\
\midrule
Base Qwen-2.5-3B & 62\%$^{\dagger}$ & 6\%$^{\dagger}$ & -- & -- \\
DeepSeek-V4-Flash & 0\% & 0\% & -- & -- \\
Gemini-2.5-Flash & $\sim$25\% & -- & -- & -- \\
\textbf{Vyasa-3B (ours)} & \textbf{66.7\%} & \textbf{25.0\%} & 64.2s & 4e/14r \\
DeepSeek-Chat (teacher) & 98\%$^{\dagger}$ & 54\%$^{\dagger}$ & $\sim$30s & -- \\
\bottomrule
\end{tabular}
\caption{Model comparison on architectural graph generation. $^{\dagger}$Re-measured 2026-07-03 under a preregistered $n{=}50$ protocol (Wilson 95\% CIs in the repository's eval reports); the earlier estimates for these rows (0\% base, $\sim$90\%/$\sim$85\% teacher) were small-$n$ and are superseded.}
\end{table}

\begin{table}[h]
\centering
\begin{tabular}{lcc}
\toprule
Component & Speed & Throughput \\
\midrule
Intent Classifier & 0.007ms & 140K/sec \\
Witness (12 checks) & 0.10ms & 123K checks/sec \\
Manifest compiler & 0.40ms & 92 files/ms \\
\textbf{Total deterministic} & \textbf{$<$2ms} & -- \\
\bottomrule
\end{tabular}
\caption{Deterministic component benchmarks}
\end{table}

\section{Discussion}
The key finding is not that a 3B model achieves 66.7\% valid JSON. It's that the 33.3\% of invalid outputs \textbf{cost zero in production}. Every invalid graph is caught by the Witness before code is generated.

Vyasa is the only system combining: (a) structured IR, (b) deterministic validation, (c) local model deployment, and (d) a training flywheel that doesn't require human evaluation.

\section{Conclusion}
The Vyasa architecture separates conception from validation from compilation. The Architectural Graph compresses specification 13:1. The Witness guarantees structural correctness. The Manifest Engine produces compilable code in under 1ms. The GRPO+Witness flywheel provides a path to improvement without human evaluation. All code is open-source at \url{https://github.com/prashantpandey-creator/vyasa-manifestation-engine}.

\begin{thebibliography}{9}
\bibitem{appycodes} Appycodes. ``We Audited 31 AI Codebases: What Actually Survives Production.'' 2026.
\bibitem{plancompiler} Chen et al. ``PlanCompiler: A Deterministic Compilation Architecture for Structured Multi-Step LLM Pipelines.'' arXiv:2604.13092, 2026.
\bibitem{compiledai} Wang et al. ``Compiled AI: Deterministic Code Generation for LLM-Based Workflow Automation.'' arXiv:2604.05150, 2026.
\bibitem{vyasa} Vyasa Manifestation Engine. GitHub: prashantpandey-creator/vyasa-manifestation-engine, 2026.
\bibitem{xvyasa} xVyasa. Hugging Face: PRASSANT/xVyasa, 2026.
\end{thebibliography}

\end{document}
